%% prototipo-tipografia.tex
%% O bloco tipográfico decidido em 17/09/2026, como protótipo do que vai
%% para a crpsp-base.sty: Luciole no corpo, NEWJUNE nos títulos, 12 pt,
%% WCAG 1.4.8 inteiro.
%%
%% Verificado em 17/09 sobre um capítulo do Guia de Apresentações
%% Acessíveis, com a mancha e o fundo decorativo reais da publicação:
%% zero erro fatal e zero erro de tagging, 4 páginas -> 5 (+25%).
%%
%% ⚠️ O documento de teste NÃO está aqui. O capítulo é conteúdo de
%% publicação e fica no workbench
%% (production/editorial/publicacoes/cartilhas/apresentacoes_acessiveis/).
%% Para repetir o teste: carregar este bloco depois do pacote da linha,
%% com a classe em 12pt.
%%
%% Como usar sobre o livros_crp_acessivel_book:
%%   \documentclass[twoside, a5paper, openright, 12pt]{book}
%%   \usepackage{livros_crp_acessivel_book}
%%   %% prototipo-tipografia.tex
%% O bloco tipográfico decidido em 17/09/2026, como protótipo do que vai
%% para a crpsp-base.sty: Luciole no corpo, NEWJUNE nos títulos, 12 pt,
%% WCAG 1.4.8 inteiro.
%%
%% Verificado em 17/09 sobre um capítulo do Guia de Apresentações
%% Acessíveis, com a mancha e o fundo decorativo reais da publicação:
%% zero erro fatal e zero erro de tagging, 4 páginas -> 5 (+25%).
%%
%% ⚠️ O documento de teste NÃO está aqui. O capítulo é conteúdo de
%% publicação e fica no workbench
%% (production/editorial/publicacoes/cartilhas/apresentacoes_acessiveis/).
%% Para repetir o teste: carregar este bloco depois do pacote da linha,
%% com a classe em 12pt.
%%
%% Como usar sobre o livros_crp_acessivel_book:
%%   \documentclass[twoside, a5paper, openright, 12pt]{book}
%%   \usepackage{livros_crp_acessivel_book}
%%   %% prototipo-tipografia.tex
%% O bloco tipográfico decidido em 17/09/2026, como protótipo do que vai
%% para a crpsp-base.sty: Luciole no corpo, NEWJUNE nos títulos, 12 pt,
%% WCAG 1.4.8 inteiro.
%%
%% Verificado em 17/09 sobre um capítulo do Guia de Apresentações
%% Acessíveis, com a mancha e o fundo decorativo reais da publicação:
%% zero erro fatal e zero erro de tagging, 4 páginas -> 5 (+25%).
%%
%% ⚠️ O documento de teste NÃO está aqui. O capítulo é conteúdo de
%% publicação e fica no workbench
%% (production/editorial/publicacoes/cartilhas/apresentacoes_acessiveis/).
%% Para repetir o teste: carregar este bloco depois do pacote da linha,
%% com a classe em 12pt.
%%
%% Como usar sobre o livros_crp_acessivel_book:
%%   \documentclass[twoside, a5paper, openright, 12pt]{book}
%%   \usepackage{livros_crp_acessivel_book}
%%   %% prototipo-tipografia.tex
%% O bloco tipográfico decidido em 17/09/2026, como protótipo do que vai
%% para a crpsp-base.sty: Luciole no corpo, NEWJUNE nos títulos, 12 pt,
%% WCAG 1.4.8 inteiro.
%%
%% Verificado em 17/09 sobre um capítulo do Guia de Apresentações
%% Acessíveis, com a mancha e o fundo decorativo reais da publicação:
%% zero erro fatal e zero erro de tagging, 4 páginas -> 5 (+25%).
%%
%% ⚠️ O documento de teste NÃO está aqui. O capítulo é conteúdo de
%% publicação e fica no workbench
%% (production/editorial/publicacoes/cartilhas/apresentacoes_acessiveis/).
%% Para repetir o teste: carregar este bloco depois do pacote da linha,
%% com a classe em 12pt.
%%
%% Como usar sobre o livros_crp_acessivel_book:
%%   \documentclass[twoside, a5paper, openright, 12pt]{book}
%%   \usepackage{livros_crp_acessivel_book}
%%   \input{prototipo-tipografia}

% ── Corpo em Luciole ────────────────────────────────────────────────
% O pacote da linha book põe Lora no corpo e NEWJUNE na sem-serifa. Só o
% corpo muda: os títulos já são \bfseries\sffamily, e continuam NEWJUNE.
% A Luciole vem do TeX Live (pacote `luciole`, quatro .ttf de texto); não
% há .sty a carregar, o uso é por fontspec.
\setmainfont{Luciole}[
    UprightFont    = Luciole-Regular.ttf,
    BoldFont       = Luciole-Bold.ttf,
    ItalicFont     = Luciole-Italic.ttf,
    BoldItalicFont = Luciole-BoldItalic.ttf,
]

% ── WCAG 1.4.8, as duas metades ─────────────────────────────────────
% A entrelinha já vem do \setasuspacing{\onehalfspacing} do pacote: 17,99 pt
% em 12 pt, porque o setspace calibra por tamanho de fonte e não pela
% entrelinha simples. O que falta é o espaço entre parágrafos, que o pacote
% deixa em 0.5\baselineskip.
\AtBeginDocument{%
  \setlength{\parskip}{1.5\baselineskip}%
  \setlength{\parindent}{0pt}%
}

% ⚠️ O \parskip global sozinho PIORA o resultado: soma-se ao afterskip do
% \@startsection (0.618\baselineskip na \section do pacote) e o título fica
% mais longe do próprio texto do que os parágrafos ficam entre si — o
% contrário do que a hierarquia pede. O critério fala de espaço ENTRE
% PARÁGRAFOS, não entre título e texto.
%
% Aqui o parágrafo logo depois de um título perde o \parskip; os demais
% mantêm. Não dá para resolver pelo afterskip: 0.618 - 1.5 é negativo, e
% afterskip negativo no \@startsection significa título embutido na linha.
\makeatletter
\let\crpsp@afterheading\@afterheading
\renewcommand{\@afterheading}{\crpsp@afterheading\vspace{-\parskip}}
\makeatother

% ── Títulos proporcionais ao corpo ──────────────────────────────────
% ⚠️ O \Huge das classes padrão SATURA: vale 24,88 pt tanto em 11 pt quanto
% em 12 pt de base. Ao subir o corpo de 10,95 para 12 pt (+9,6%), o título de
% capítulo -- que é `title-decls = \Huge` no pacote book -- NÃO cresce, e a
% razão título/corpo cai de 2,27x para 2,07x. O número do capítulo (\huge)
% faz o contrário: cresce 20%, mais que o corpo. Ou seja, subir o corpo
% comprime a hierarquia em silêncio.
%
% Os dois voltam à proporção do desenho original, multiplicando os valores de
% 11 pt por 12/10,95:
%     \Huge  24,88 -> 27,27        \huge  20,74 -> 22,73
%
% A entrelinha dos títulos fica em 1,15x o corpo da fonte em vez de herdar a
% entrelinha de 1,5 do texto: num título de três linhas o 1,5 solta demais.
% Isto NÃO fere o WCAG 1.4.8, cujo critério vale para blocos de texto.
\renewcommand{\Huge}{\fontsize{27.27}{31.36}\selectfont}
\renewcommand{\huge}{\fontsize{22.73}{26.14}\selectfont}

% ── Alinhamento à esquerda, sem hifenação ───────────────────────────
% "Texto não justificado" é um dos cinco itens do próprio WCAG 1.4.8, junto
% com a largura de linha e a entrelinha. O texto justificado abre buracos
% irregulares entre palavras, que atrapalham quem tem dislexia e quem lê com
% ampliação. A hifenação sai à parte: quebra de palavra obriga a remontar o
% termo, e o alinhamento à esquerda não precisa dela para fechar a linha.
% \exhyphenpenalty cuida dos hifens já presentes no texto ("cão-guia").
%
% ⚠️ O \RaggedRight do ragged2e sozinho NÃO basta. O esticamento padrão dele
% à direita é finito (2em): serrilha menos a margem, mas conta com a
% hifenação para fechar as linhas difíceis. Sem ela as linhas transbordam --
% medido no capítulo de teste, 8 a 9 linhas por versão passavam da mancha,
% uma delas em 24 pt, 8 mm de texto dentro da margem. Com esticamento
% infinito não há transbordo possível por espacejamento: a linha termina
% mais cedo, que é o que o alinhamento à esquerda quer dizer.
\setlength{\RaggedRightRightskip}{0pt plus 1fil}
\AtBeginDocument{%
  \RaggedRight
  \hyphenpenalty=10000
  \exhyphenpenalty=10000
}

\endinput


% ── Corpo em Luciole ────────────────────────────────────────────────
% O pacote da linha book põe Lora no corpo e NEWJUNE na sem-serifa. Só o
% corpo muda: os títulos já são \bfseries\sffamily, e continuam NEWJUNE.
% A Luciole vem do TeX Live (pacote `luciole`, quatro .ttf de texto); não
% há .sty a carregar, o uso é por fontspec.
\setmainfont{Luciole}[
    UprightFont    = Luciole-Regular.ttf,
    BoldFont       = Luciole-Bold.ttf,
    ItalicFont     = Luciole-Italic.ttf,
    BoldItalicFont = Luciole-BoldItalic.ttf,
]

% ── WCAG 1.4.8, as duas metades ─────────────────────────────────────
% A entrelinha já vem do \setasuspacing{\onehalfspacing} do pacote: 17,99 pt
% em 12 pt, porque o setspace calibra por tamanho de fonte e não pela
% entrelinha simples. O que falta é o espaço entre parágrafos, que o pacote
% deixa em 0.5\baselineskip.
\AtBeginDocument{%
  \setlength{\parskip}{1.5\baselineskip}%
  \setlength{\parindent}{0pt}%
}

% ⚠️ O \parskip global sozinho PIORA o resultado: soma-se ao afterskip do
% \@startsection (0.618\baselineskip na \section do pacote) e o título fica
% mais longe do próprio texto do que os parágrafos ficam entre si — o
% contrário do que a hierarquia pede. O critério fala de espaço ENTRE
% PARÁGRAFOS, não entre título e texto.
%
% Aqui o parágrafo logo depois de um título perde o \parskip; os demais
% mantêm. Não dá para resolver pelo afterskip: 0.618 - 1.5 é negativo, e
% afterskip negativo no \@startsection significa título embutido na linha.
\makeatletter
\let\crpsp@afterheading\@afterheading
\renewcommand{\@afterheading}{\crpsp@afterheading\vspace{-\parskip}}
\makeatother

% ── Títulos proporcionais ao corpo ──────────────────────────────────
% ⚠️ O \Huge das classes padrão SATURA: vale 24,88 pt tanto em 11 pt quanto
% em 12 pt de base. Ao subir o corpo de 10,95 para 12 pt (+9,6%), o título de
% capítulo -- que é `title-decls = \Huge` no pacote book -- NÃO cresce, e a
% razão título/corpo cai de 2,27x para 2,07x. O número do capítulo (\huge)
% faz o contrário: cresce 20%, mais que o corpo. Ou seja, subir o corpo
% comprime a hierarquia em silêncio.
%
% Os dois voltam à proporção do desenho original, multiplicando os valores de
% 11 pt por 12/10,95:
%     \Huge  24,88 -> 27,27        \huge  20,74 -> 22,73
%
% A entrelinha dos títulos fica em 1,15x o corpo da fonte em vez de herdar a
% entrelinha de 1,5 do texto: num título de três linhas o 1,5 solta demais.
% Isto NÃO fere o WCAG 1.4.8, cujo critério vale para blocos de texto.
\renewcommand{\Huge}{\fontsize{27.27}{31.36}\selectfont}
\renewcommand{\huge}{\fontsize{22.73}{26.14}\selectfont}

% ── Alinhamento à esquerda, sem hifenação ───────────────────────────
% "Texto não justificado" é um dos cinco itens do próprio WCAG 1.4.8, junto
% com a largura de linha e a entrelinha. O texto justificado abre buracos
% irregulares entre palavras, que atrapalham quem tem dislexia e quem lê com
% ampliação. A hifenação sai à parte: quebra de palavra obriga a remontar o
% termo, e o alinhamento à esquerda não precisa dela para fechar a linha.
% \exhyphenpenalty cuida dos hifens já presentes no texto ("cão-guia").
%
% ⚠️ O \RaggedRight do ragged2e sozinho NÃO basta. O esticamento padrão dele
% à direita é finito (2em): serrilha menos a margem, mas conta com a
% hifenação para fechar as linhas difíceis. Sem ela as linhas transbordam --
% medido no capítulo de teste, 8 a 9 linhas por versão passavam da mancha,
% uma delas em 24 pt, 8 mm de texto dentro da margem. Com esticamento
% infinito não há transbordo possível por espacejamento: a linha termina
% mais cedo, que é o que o alinhamento à esquerda quer dizer.
\setlength{\RaggedRightRightskip}{0pt plus 1fil}
\AtBeginDocument{%
  \RaggedRight
  \hyphenpenalty=10000
  \exhyphenpenalty=10000
}

\endinput


% ── Corpo em Luciole ────────────────────────────────────────────────
% O pacote da linha book põe Lora no corpo e NEWJUNE na sem-serifa. Só o
% corpo muda: os títulos já são \bfseries\sffamily, e continuam NEWJUNE.
% A Luciole vem do TeX Live (pacote `luciole`, quatro .ttf de texto); não
% há .sty a carregar, o uso é por fontspec.
\setmainfont{Luciole}[
    UprightFont    = Luciole-Regular.ttf,
    BoldFont       = Luciole-Bold.ttf,
    ItalicFont     = Luciole-Italic.ttf,
    BoldItalicFont = Luciole-BoldItalic.ttf,
]

% ── WCAG 1.4.8, as duas metades ─────────────────────────────────────
% A entrelinha já vem do \setasuspacing{\onehalfspacing} do pacote: 17,99 pt
% em 12 pt, porque o setspace calibra por tamanho de fonte e não pela
% entrelinha simples. O que falta é o espaço entre parágrafos, que o pacote
% deixa em 0.5\baselineskip.
\AtBeginDocument{%
  \setlength{\parskip}{1.5\baselineskip}%
  \setlength{\parindent}{0pt}%
}

% ⚠️ O \parskip global sozinho PIORA o resultado: soma-se ao afterskip do
% \@startsection (0.618\baselineskip na \section do pacote) e o título fica
% mais longe do próprio texto do que os parágrafos ficam entre si — o
% contrário do que a hierarquia pede. O critério fala de espaço ENTRE
% PARÁGRAFOS, não entre título e texto.
%
% Aqui o parágrafo logo depois de um título perde o \parskip; os demais
% mantêm. Não dá para resolver pelo afterskip: 0.618 - 1.5 é negativo, e
% afterskip negativo no \@startsection significa título embutido na linha.
\makeatletter
\let\crpsp@afterheading\@afterheading
\renewcommand{\@afterheading}{\crpsp@afterheading\vspace{-\parskip}}
\makeatother

% ── Títulos proporcionais ao corpo ──────────────────────────────────
% ⚠️ O \Huge das classes padrão SATURA: vale 24,88 pt tanto em 11 pt quanto
% em 12 pt de base. Ao subir o corpo de 10,95 para 12 pt (+9,6%), o título de
% capítulo -- que é `title-decls = \Huge` no pacote book -- NÃO cresce, e a
% razão título/corpo cai de 2,27x para 2,07x. O número do capítulo (\huge)
% faz o contrário: cresce 20%, mais que o corpo. Ou seja, subir o corpo
% comprime a hierarquia em silêncio.
%
% Os dois voltam à proporção do desenho original, multiplicando os valores de
% 11 pt por 12/10,95:
%     \Huge  24,88 -> 27,27        \huge  20,74 -> 22,73
%
% A entrelinha dos títulos fica em 1,15x o corpo da fonte em vez de herdar a
% entrelinha de 1,5 do texto: num título de três linhas o 1,5 solta demais.
% Isto NÃO fere o WCAG 1.4.8, cujo critério vale para blocos de texto.
\renewcommand{\Huge}{\fontsize{27.27}{31.36}\selectfont}
\renewcommand{\huge}{\fontsize{22.73}{26.14}\selectfont}

% ── Alinhamento à esquerda, sem hifenação ───────────────────────────
% "Texto não justificado" é um dos cinco itens do próprio WCAG 1.4.8, junto
% com a largura de linha e a entrelinha. O texto justificado abre buracos
% irregulares entre palavras, que atrapalham quem tem dislexia e quem lê com
% ampliação. A hifenação sai à parte: quebra de palavra obriga a remontar o
% termo, e o alinhamento à esquerda não precisa dela para fechar a linha.
% \exhyphenpenalty cuida dos hifens já presentes no texto ("cão-guia").
%
% ⚠️ O \RaggedRight do ragged2e sozinho NÃO basta. O esticamento padrão dele
% à direita é finito (2em): serrilha menos a margem, mas conta com a
% hifenação para fechar as linhas difíceis. Sem ela as linhas transbordam --
% medido no capítulo de teste, 8 a 9 linhas por versão passavam da mancha,
% uma delas em 24 pt, 8 mm de texto dentro da margem. Com esticamento
% infinito não há transbordo possível por espacejamento: a linha termina
% mais cedo, que é o que o alinhamento à esquerda quer dizer.
\setlength{\RaggedRightRightskip}{0pt plus 1fil}
\AtBeginDocument{%
  \RaggedRight
  \hyphenpenalty=10000
  \exhyphenpenalty=10000
}

\endinput


% ── Corpo em Luciole ────────────────────────────────────────────────
% O pacote da linha book põe Lora no corpo e NEWJUNE na sem-serifa. Só o
% corpo muda: os títulos já são \bfseries\sffamily, e continuam NEWJUNE.
% A Luciole vem do TeX Live (pacote `luciole`, quatro .ttf de texto); não
% há .sty a carregar, o uso é por fontspec.
\setmainfont{Luciole}[
    UprightFont    = Luciole-Regular.ttf,
    BoldFont       = Luciole-Bold.ttf,
    ItalicFont     = Luciole-Italic.ttf,
    BoldItalicFont = Luciole-BoldItalic.ttf,
]

% ── WCAG 1.4.8, as duas metades ─────────────────────────────────────
% A entrelinha já vem do \setasuspacing{\onehalfspacing} do pacote: 17,99 pt
% em 12 pt, porque o setspace calibra por tamanho de fonte e não pela
% entrelinha simples. O que falta é o espaço entre parágrafos, que o pacote
% deixa em 0.5\baselineskip.
\AtBeginDocument{%
  \setlength{\parskip}{1.5\baselineskip}%
  \setlength{\parindent}{0pt}%
}

% ⚠️ O \parskip global sozinho PIORA o resultado: soma-se ao afterskip do
% \@startsection (0.618\baselineskip na \section do pacote) e o título fica
% mais longe do próprio texto do que os parágrafos ficam entre si — o
% contrário do que a hierarquia pede. O critério fala de espaço ENTRE
% PARÁGRAFOS, não entre título e texto.
%
% Aqui o parágrafo logo depois de um título perde o \parskip; os demais
% mantêm. Não dá para resolver pelo afterskip: 0.618 - 1.5 é negativo, e
% afterskip negativo no \@startsection significa título embutido na linha.
\makeatletter
\let\crpsp@afterheading\@afterheading
\renewcommand{\@afterheading}{\crpsp@afterheading\vspace{-\parskip}}
\makeatother

% ── Títulos proporcionais ao corpo ──────────────────────────────────
% ⚠️ O \Huge das classes padrão SATURA: vale 24,88 pt tanto em 11 pt quanto
% em 12 pt de base. Ao subir o corpo de 10,95 para 12 pt (+9,6%), o título de
% capítulo -- que é `title-decls = \Huge` no pacote book -- NÃO cresce, e a
% razão título/corpo cai de 2,27x para 2,07x. O número do capítulo (\huge)
% faz o contrário: cresce 20%, mais que o corpo. Ou seja, subir o corpo
% comprime a hierarquia em silêncio.
%
% Os dois voltam à proporção do desenho original, multiplicando os valores de
% 11 pt por 12/10,95:
%     \Huge  24,88 -> 27,27        \huge  20,74 -> 22,73
%
% A entrelinha dos títulos fica em 1,15x o corpo da fonte em vez de herdar a
% entrelinha de 1,5 do texto: num título de três linhas o 1,5 solta demais.
% Isto NÃO fere o WCAG 1.4.8, cujo critério vale para blocos de texto.
\renewcommand{\Huge}{\fontsize{27.27}{31.36}\selectfont}
\renewcommand{\huge}{\fontsize{22.73}{26.14}\selectfont}

% ── Alinhamento à esquerda, sem hifenação ───────────────────────────
% "Texto não justificado" é um dos cinco itens do próprio WCAG 1.4.8, junto
% com a largura de linha e a entrelinha. O texto justificado abre buracos
% irregulares entre palavras, que atrapalham quem tem dislexia e quem lê com
% ampliação. A hifenação sai à parte: quebra de palavra obriga a remontar o
% termo, e o alinhamento à esquerda não precisa dela para fechar a linha.
% \exhyphenpenalty cuida dos hifens já presentes no texto ("cão-guia").
%
% ⚠️ O \RaggedRight do ragged2e sozinho NÃO basta. O esticamento padrão dele
% à direita é finito (2em): serrilha menos a margem, mas conta com a
% hifenação para fechar as linhas difíceis. Sem ela as linhas transbordam --
% medido no capítulo de teste, 8 a 9 linhas por versão passavam da mancha,
% uma delas em 24 pt, 8 mm de texto dentro da margem. Com esticamento
% infinito não há transbordo possível por espacejamento: a linha termina
% mais cedo, que é o que o alinhamento à esquerda quer dizer.
\setlength{\RaggedRightRightskip}{0pt plus 1fil}
\AtBeginDocument{%
  \RaggedRight
  \hyphenpenalty=10000
  \exhyphenpenalty=10000
}

\endinput
